\answersfor{ch:scoringgains}

\begin{solution}{ex:scoringgains:1}
\emph{Invert.} $ad - bc = 1 \times 1 - 0.5 \times 0.5 = 0.75$, so
\[
\Cov\inv = \frac{1}{0.75}\begin{pmatrix} 1 & -0.5 \\ -0.5 & 1\end{pmatrix} = \begin{pmatrix} 1.333 & -0.667 \\ -0.667 & 1.333\end{pmatrix}.
\]
\emph{Apply to $\vect{t} = (1, 1)$.} $\Cov\inv\vect{t} = (1.333 - 0.667,\ -0.667 + 1.333) = (0.667, 0.667)$;
$\vect{t}\T\Cov\inv\vect{t} = 0.667 + 0.667 = 1.333$.

\emph{Filter and SNR.} $\vect{w}^\star = (0.667, 0.667) / 1.333 = (0.5, 0.5)$; SNR $= \sqrt{1.333} = 1.155$.

\emph{Naive $\vect{w} = \vect{t}$.} $\vect{w}\T\vect{t} = 1 + 1 = 2$; $\vect{w}\T\Cov\vect{w} = 1 + 0.5 + 0.5 + 1 = 3$; SNR $=
2 / \sqrt{3} = 1.155$. The same. Here $\vect{t}$ is an eigenvector of $\Cov$ ($\Cov\vect{t} = (1.5, 1.5) = 1.5\,\vect{t}$),
so $\Cov\inv\vect{t} = \vect{t} / 1.5$ points the same way as $\vect{t}$: the matched filter is the naive filter, rescaled.
Matching only helps when the noise's correlations point away from the signal, as in \cref{we:scoringgains:mf2}.
\end{solution}

\begin{solution}{ex:scoringgains:2}
$b_{11} = 0.25 \times 0 + 0.5 \times 5.2 + 0.25 \times 0 = 2.6$, score $2.6^2 = 6.76$. The blurred ranking becomes 18.06
(pixel 7), 16 (6), 7.56 (8), 6.76 (11), 6.25 (5): the target is fourth and pixel 5 drops to fifth. Both objects are now
caught - the big one through pixels 6--8 - so objects caught returns to 2. Four tenths of a unit of residual decided whether
the one-pixel target survived blurring; a filter that puts a small target this close to the cut-off loses many of them on
a real scene.
\end{solution}

\begin{solution}{ex:scoringgains:3}
New PCA-RX scores: P1 11, P2 1, P3 2, P4 10, P5 3. Sorted: P2 (1), P3 (2), P5 (3), P4 (10), P1 (11), so the PCA-RX ranks
are P2 $1/5$, P3 $2/5$, P5 $3/5$, P4 $4/5$, P1 $5/5$. RM ranks are unchanged: P1 $4/5$, P4 $2/5$, P2 $5/5$.
\begin{align*}
\text{P1:}\quad & \text{mean } (4/5 + 5/5)/2 = 9/10 = 0.9, \quad \text{max } 1.0\\
\text{P4:}\quad & \text{mean } (2/5 + 4/5)/2 = 6/10 = 0.6, \quad \text{max } 0.8\\
\text{P2:}\quad & \text{mean } (5/5 + 1/5)/2 = 6/10 = 0.6, \quad \text{max } 1.0
\end{align*}
Rank mean puts P1 first (0.9). Rank max ties P1 with P2, RM's false alarm, at 1.0.
\end{solution}

\begin{solution}{ex:scoringgains:4}
\begin{align*}
\mathcal{N}(2; 0, 1) &= 0.3989 \times e^{-2} = 0.3989 \times 0.1353 = 0.05399\\
\mathcal{N}(2; 5, 4) &= 0.1995 \times e^{-9/8} = 0.1995 \times 0.3247 = 0.06476\\
\pi_1\mathcal{N}_1 &= 0.8 \times 0.05399 = 0.04319, \qquad \pi_2\mathcal{N}_2 = 0.2 \times 0.06476 = 0.01295\\
r_1 &= 0.04319 / (0.04319 + 0.01295) = 0.04319 / 0.05614 = 0.769, \qquad r_2 = 0.231
\end{align*}
Squared distances: field $(2 - 0)^2/1 = 4$; roofs $(2 - 5)^2/4 = 9/4 = 2.25$. Weighted score $0.769 \times 4 + 0.231 \times
2.25 = 3.077 + 0.519 = 3.60$. (Script: 3.596.)
\end{solution}

\begin{solution}{ex:scoringgains:5}
Eight built-up pixels: six roofs at $\pm 2$ and two targets at $\pm 6$.
\begin{align*}
\text{sum of squares} &= 6 \times 4 + 2 \times 36 = 24 + 72 = 96\\
\text{variance} &= 96 / 8 = 12\\
\text{target score} &= 36 / 12 = 3.0
\end{align*}
Still below the field edges' 3.77: the targets lose. Diluting them with more ordinary roofs helps, but not enough; on set4 one
type holds 60\% of the target pixels, far more concentrated than this.
\end{solution}

\begin{solution}{ex:scoringgains:6}
Only $N = 50$ and $N = 200$:
\begin{itemize}
\item set1 held out: $0.1393$ against $0.1360$, picks 50, reports 0.223.
\item set2 held out: $0.1663$ against $0.1647$, picks 50, reports 0.142.
\item set3 held out: $0.1857$ against $0.1850$, picks 50, reports 0.084.
\item set4 held out: $0.1497$ against $0.1503$, picks 200, reports 0.185.
\end{itemize}
Headline: $(0.223 + 0.142 + 0.084 + 0.185)/4 = 0.634/4 = 0.1585$, a gain of $0.1585 - 0.1695 = -0.011$. Different candidates,
different picks - and still below FT04. The choice of $N$ is unstable; the verdict is not.
\end{solution}

\begin{solution}{ex:scoringgains:7}
With $a = 1$, $b = 4$ in \eqref{eq:scoringgains:dilate}: $1 \times 4 + 4 \times (1 + 4) + 4 = 4 + 20 + 4 = 28$ pixels (SciPy
agrees). Growth $28 / 4 = 7$, so a top 1\% made of such lines would gate about $1\% \times 7 = 7\%$ of the scene - between
the $2 \times 2$ clumps' 6\% and isolated pixels' 13\%.
\end{solution}
